\documentclass[11pt]{article}
\usepackage[margin=1in]{geometry}
\usepackage{amsmath,amssymb}
\usepackage{booktabs}
\usepackage{algorithm}
\usepackage{algpseudocode}
\usepackage{hyperref}
\usepackage{xcolor}

\title{ARG Traversals and Tree-Shuttle Multiply\\from Threading Instructions}
\author{Technical Note}
\date{March 2026}

\begin{document}
\maketitle

\begin{abstract}
We describe algorithms for direct ARG traversal and genotype-matrix
multiplication from \textsc{Threads} threading instructions.  The key
contribution is a family of \emph{visit} algorithms that reconstruct local
trees via union-find at segment boundaries, avoiding materialization of the
full ARG or genotype matrix.  Three specializations---full member tracking,
count-only, and height-only---yield efficient implementations of clade
enumeration, allele frequency spectrum, total branch volume, association
testing, and mutation generation.  We also describe the \emph{tree-shuttle}
multiply algorithm, which propagates cumulative sums through the
interval-tree structure of threading instructions for implicit
genotype-matrix--vector products in $O(n \cdot n_\text{intervals} +
n_\text{mismatches})$ time per call.
\end{abstract}

%% ========================================================================
\section{Preliminaries}

\subsection{Threading Instructions}

The \textsc{Threads} software infers an Ancestral Recombination Graph (ARG)
by threading each haplotype through a set of reference haplotypes.  For each
of $n$ haplotype samples, the result is a sequence of \emph{segments}
$\{(\text{start}_k, \text{target}_k, \text{tmrca}_k)\}$ and a list of
\emph{mismatch} site indices.  In each segment, the sample copies the
genotype of its target, except at mismatch sites where the allele is
flipped.  The coalescence time $\text{tmrca}_k$ records the height at which
the sample and its target share a most recent common ancestor in segment
$k$.

Samples are ordered $0, 1, \ldots, n{-}1$.  Sample~0 is the reference
haplotype (no target; genotype determined entirely by mismatches from
all-zero).  For $i > 0$, every segment's target is a sample with index $<
i$, ensuring a well-defined dependency order.

\subsection{Segment Boundaries and Local Trees}

A \emph{boundary position} is any genomic position where at least one
sample's segment changes.  Let $B = \{b_0, b_1, \ldots, b_{K-1}\}$ be the
sorted set of all segment start positions across all samples.  Between
consecutive boundaries $[b_j, b_{j+1})$, the coalescence relationships are
constant: every sample has the same (target, tmrca) pair throughout the
interval.  We call $[b_j, b_{j+1})$ an \emph{interval}.

At each interval, the threading instructions define a forest of coalescent
trees: sample $i$ coalesces with its target at height $\text{tmrca}_i$.
Reconstructing these trees is the basis for all visit algorithms.

%% ========================================================================
\section{Union-Find Tree Reconstruction}

All visit algorithms share a common structure: iterate over intervals,
reconstruct the local tree via union-find, and invoke a callback for each
clade or branch.  The key variation is the \emph{level of detail} tracked by
the union-find:

\begin{enumerate}
\item \textbf{Full member tracking}: Each component stores a sorted list of
all descendant leaf indices.  Required for \texttt{visit\_clades},
\texttt{visit\_branches}, and \texttt{association\_diploid}.
\item \textbf{Count-only}: Each component stores only its size (number of
descendants).  Sufficient for \texttt{allele\_frequency\_spectrum}.
\item \textbf{Height-only}: Each component stores only its coalescence
height.  Sufficient for \texttt{total\_volume}.
\end{enumerate}

\subsection{Data Structures}

For full member tracking, the union-find stores:
\begin{itemize}
\item $\texttt{par}[n]$: parent pointers with path compression
\item $\texttt{rnk}[n]$: rank for union-by-rank
\item $\texttt{members}[n]$: vector of descendant indices per root
\item $\texttt{height}[n]$: coalescence height per root
\end{itemize}

For count-only and height-only variants, $\texttt{members}$ is replaced by
$\texttt{size}[n]$ (integer count) or omitted entirely.

\subsection{Identical-Tree Optimization}

Consecutive intervals often share the same tree topology (same target and
tmrca for every sample).  Rather than rebuilding the union-find redundantly,
we detect this by comparing the current interval's (target, tmrca) vectors
against the previous interval's.  When the topology is unchanged:

\begin{itemize}
\item \textbf{visit\_clades / visit\_branches}: Pending clade/branch records
are extended to span the new interval rather than emitted and recreated.
\item \textbf{total\_volume}: The per-interval volume contribution is
reused: $\text{vol} \mathrel{+}= \text{prev\_vol} \times \text{span}$.
\item \textbf{AFS}: The per-bin volume contribution is reused similarly.
\end{itemize}

This optimization is particularly effective in regions of low recombination,
where many consecutive intervals share the same tree.

%% ========================================================================
\section{Visit Clades}

\begin{algorithm}[h]
\caption{\texttt{visit\_clades}: Enumerate all clades in the ARG}
\begin{algorithmic}[1]
\State Collect boundary positions $B$ from all samples' segment starts
\State Emit leaf clades: for each $i \in [0, n)$, callback$(\{i\}, 0, \text{start}, \text{end})$
\For{each interval $[b_j, b_{j+1})$}
    \State Look up $(\text{target}_i, \text{tmrca}_i)$ for each sample $i$ via binary search on segment starts
    \If{topology unchanged from previous interval}
        \State \textbf{continue} \Comment{extend pending clades}
    \EndIf
    \State Flush pending clades with callback
    \State Reset union-find: each sample is its own component
    \State Collect edges: $E = \{(i, \text{target}_i, \text{tmrca}_i) : i \ge 1, \text{target}_i \ge 0\}$
    \State Sort $E$ by tmrca (ascending)
    \For{each edge $(c, p, h)$ in sorted $E$}
        \State $r_c \gets \text{find}(c)$, $r_p \gets \text{find}(p)$
        \If{$r_c \ne r_p$}
            \State $\text{root} \gets \text{unite}(r_c, r_p)$
            \State $\texttt{height}[\text{root}] \gets h$
            \State $\texttt{members}[\text{root}] \gets \texttt{members}[r_c] \cup \texttt{members}[r_p]$
            \State Add pending clade: $(\texttt{members}[\text{root}], h, b_j)$
        \EndIf
    \EndFor
\EndFor
\State Flush remaining pending clades
\end{algorithmic}
\end{algorithm}

\textbf{Callback signature:}
$\text{callback}(\texttt{descendants}, \texttt{height},
\texttt{start\_bp}, \texttt{end\_bp})$, where \texttt{descendants} is the
sorted list of leaf indices in the clade, \texttt{height} is the
coalescence time, and the interval $[\texttt{start\_bp},
\texttt{end\_bp})$ is the genomic span.

\textbf{Complexity:}  $O(|B| \cdot n \log n)$ time ($n \log n$ per interval
for sorting edges and union-find operations), $O(n)$ space for the
union-find.

%% ========================================================================
\section{Visit Branches}

\texttt{visit\_branches} follows the same structure as
\texttt{visit\_clades} but emits \emph{branches} (edges between tree nodes)
rather than clades.

\begin{algorithm}[h]
\caption{\texttt{visit\_branches}: Enumerate all branches in the ARG}
\begin{algorithmic}[1]
\State (Same boundary collection and topology detection as \texttt{visit\_clades})
\For{each edge $(c, p, h)$ uniting components $r_c, r_p$}
    \If{$\texttt{height}[r_c] < h$}
        \State Pending branch: $(\texttt{members}[r_c], \texttt{height}[r_c], h, b_j)$
    \EndIf
    \If{$\texttt{height}[r_p] < h$}
        \State Pending branch: $(\texttt{members}[r_p], \texttt{height}[r_p], h, b_j)$
    \EndIf
    \State $\text{root} \gets \text{unite}(r_c, r_p)$; $\texttt{height}[\text{root}] \gets h$
\EndFor
\end{algorithmic}
\end{algorithm}

\textbf{Callback signature:}
$\text{callback}(\texttt{descendants}, \texttt{child\_height},
\texttt{parent\_height}, \texttt{start\_bp}, \texttt{end\_bp})$.

A branch represents a tree edge from a node at \texttt{child\_height} to a
node at \texttt{parent\_height}, subtending the samples in
\texttt{descendants}.  Each coalescence event produces up to two branches
(one for each merging component whose height is below the coalescence
time).

%% ========================================================================
\section{Total Volume}

\texttt{total\_volume()} computes the total branch length of the ARG:
\[
V = \sum_{\text{branches}} (\text{parent\_height} - \text{child\_height})
    \times (\text{end\_bp} - \text{start\_bp}).
\]

This uses a \emph{height-only} union-find that tracks only the coalescence
height of each component, without member lists.

\begin{algorithm}[h]
\caption{\texttt{total\_volume}: Total ARG branch volume}
\begin{algorithmic}[1]
\State $V \gets 0$
\For{each interval $[b_j, b_{j+1})$ with span $s = b_{j+1} - b_j$}
    \If{topology unchanged}
        \State $V \mathrel{+}= v_\text{prev} \times s$ \Comment{reuse cached contribution}
        \State \textbf{continue}
    \EndIf
    \State Reset height-only UF; $v_\text{prev} \gets 0$
    \State Sort edges by height
    \For{each edge $(c, p, h)$}
        \State $r_c \gets \text{find}(c)$, $r_p \gets \text{find}(p)$
        \State $v_\text{prev} \mathrel{+}= (h - \texttt{height}[r_c]) + (h - \texttt{height}[r_p])$
        \State Unite $r_c, r_p$; set $\texttt{height}[\text{root}] \gets h$
    \EndFor
    \State $V \mathrel{+}= v_\text{prev} \times s$
\EndFor
\State \Return $V$
\end{algorithmic}
\end{algorithm}

\textbf{Complexity:} $O(|B| \cdot n \log n)$ time, $O(n)$ space.  The
identical-tree optimization avoids rebuilding the union-find when the tree
is unchanged, reducing the effective cost to $O(|B_\text{unique}| \cdot n
\log n)$ where $|B_\text{unique}|$ is the number of \emph{distinct}
topologies.

%% ========================================================================
\section{Allele Frequency Spectrum}

The branch-length allele frequency spectrum (AFS) bins total branch volume
by the number of descendants:
\[
\text{afs}[k] = \sum_{\substack{\text{branches with} \\ k \text{ descendants}}}
    (\text{parent\_height} - \text{child\_height})
    \times (\text{end\_bp} - \text{start\_bp}).
\]

This uses a \emph{count-only} union-find that tracks the size (descendant
count) and height of each component, but not member lists.

\begin{algorithm}[h]
\caption{\texttt{allele\_frequency\_spectrum}: Branch-length AFS}
\begin{algorithmic}[1]
\State Initialize $\text{afs}[0 \ldots n] \gets 0$
\For{each interval $[b_j, b_{j+1})$ with span $s$}
    \If{topology unchanged}
        \For{$k = 0$ to $n$}
            \State $\text{afs}[k] \mathrel{+}= \text{afs\_prev}[k] \times s$
        \EndFor
        \State \textbf{continue}
    \EndIf
    \State Reset count-only UF: $\texttt{size}[i] \gets 1$, $\texttt{height}[i] \gets 0$
    \State Clear $\text{afs\_prev}[0 \ldots n]$
    \State Sort edges by height
    \For{each edge $(c, p, h)$ uniting $r_c, r_p$}
        \If{$\texttt{height}[r_c] < h$}
            \State $\text{afs\_prev}[\texttt{size}[r_c]] \mathrel{+}= h - \texttt{height}[r_c]$
        \EndIf
        \If{$\texttt{height}[r_p] < h$}
            \State $\text{afs\_prev}[\texttt{size}[r_p]] \mathrel{+}= h - \texttt{height}[r_p]$
        \EndIf
        \State $\text{root} \gets \text{unite}(r_c, r_p)$
        \State $\texttt{size}[\text{root}] \gets \texttt{size}[r_c] + \texttt{size}[r_p]$
        \State $\texttt{height}[\text{root}] \gets h$
    \EndFor
    \For{$k = 0$ to $n$}
        \State $\text{afs}[k] \mathrel{+}= \text{afs\_prev}[k] \times s$
    \EndFor
\EndFor
\State \Return $\text{afs}$
\end{algorithmic}
\end{algorithm}

\textbf{Complexity:} $O(|B| \cdot n \log n)$ time, $O(n)$ space.

\textbf{Comparison with arg-needle-lib:}  The existing \texttt{arg-needle-lib}
implementation computes the AFS via \texttt{bitset\_volume\_map}, which
maintains a bitset of size $n$ per tree node to track descendants.  This
requires $O(n^2 / 64)$ space and $O(n^2 / 64)$ bitwise-OR operations per
coalescence, totaling $O(|B| \cdot n^2)$.  The count-only union-find
reduces this to $O(|B| \cdot n \log n)$---a 50--70$\times$ speedup in
practice for $n \ge 1000$.

%% ========================================================================
\section{Association Testing}

\texttt{association\_diploid} performs chi-square association testing of
branch-defined clades against a diploid phenotype vector, without
materializing the genotype matrix.

\begin{algorithm}[h]
\caption{\texttt{association\_diploid}: Clade-based association}
\begin{algorithmic}[1]
\State Precompute phenotype statistics: $\bar{y}$, $\text{Var}(y)$
\State Compute $\chi^2$ threshold from $p$-value threshold
\For{each clade $(\texttt{desc}, h, \text{start}, \text{end})$ via \texttt{visit\_clades}}
    \If{$|\texttt{desc}| < 2$ or $|\texttt{desc}| \ge n - 1$}
        \State \textbf{continue} \Comment{skip trivial clades}
    \EndIf
    \State Compute diploid dosage: $d_j = |\{h \in \texttt{desc} : \lfloor h/2 \rfloor = j\}|$ for each individual $j$
    \State $\bar{d} = \sum d_j / n_\text{dip}$, \; $\text{Var}(d) = \sum d_j^2 / n_\text{dip} - \bar{d}^2$
    \State $\text{Cov}(d, y) = \sum d_j \, y_j / n_\text{dip} - \bar{d}\,\bar{y}$
    \State $r^2 = \text{Cov}(d, y)^2 / (\text{Var}(d) \cdot \text{Var}(y))$
    \State $\chi^2 = r^2 \cdot n_\text{dip}$
    \If{$\chi^2 \ge \chi^2_\text{thresh}$}
        \State Compute $p$-value via $\text{erfc}(\sqrt{\chi^2 / 2})$
        \State Emit $(\text{midpoint}, \chi^2, p, |\texttt{desc}|)$
    \EndIf
\EndFor
\end{algorithmic}
\end{algorithm}

\textbf{Dosage computation:}  For each haploid descendant $h$, the diploid
individual is $\lfloor h/2 \rfloor$.  The dosage $d_j \in \{0, 1, 2\}$
counts how many of an individual's two haplotypes are in the clade.

\textbf{Complexity:}  $O(C \cdot n_\text{dip})$ where $C$ is the number of
non-trivial clades.  The visit\_clades traversal dominates setup cost at
$O(|B| \cdot n \log n)$.

%% ========================================================================
\section{Mutation Generation}

\texttt{generate\_mutations} samples mutations on ARG branches
proportionally to branch volume, producing a genotype matrix of
derived alleles.

\begin{algorithm}[h]
\caption{\texttt{generate\_mutations}: Poisson mutation sampling}
\begin{algorithmic}[1]
\State Initialize RNG with seed
\For{each branch $(\texttt{desc}, h_c, h_p, \text{start}, \text{end})$ via \texttt{visit\_branches}}
    \State $\lambda = \mu \cdot (h_p - h_c) \cdot (\text{end} - \text{start})$
    \State $k \sim \text{Poisson}(\lambda)$
    \For{$j = 1$ to $k$}
        \State $\text{pos} \sim \text{Uniform}(\text{start}, \text{end})$
        \State Set genotype: 1 for samples in \texttt{desc}, 0 otherwise
    \EndFor
\EndFor
\State \Return positions, genotypes (flat $n_\text{mut} \times n$ matrix)
\end{algorithmic}
\end{algorithm}

\textbf{Complexity:}  $O(B_\text{branches} + n_\text{mut} \cdot n)$ where
$B_\text{branches}$ is the cost of \texttt{visit\_branches} and
$n_\text{mut}$ is the total number of sampled mutations.  The expected
number of mutations is $\mu \cdot V$ where $V$ is the total volume.

%% ========================================================================
\section{Tree-Shuttle Multiply}

The tree-shuttle multiply computes implicit genotype-matrix--vector products
$\mathbf{G}^\top \mathbf{x}$ (right multiply) and $\mathbf{G}\mathbf{x}$
(left multiply) by propagating cumulative sums through the dependency
structure of threading instructions, without materializing the genotype
matrix.  Both directions use $O(n)$ per-call memory and achieve sublinear
scaling in practice.

\subsection{Notation}

Throughout this section we use the following symbols:
\begin{itemize}
\item $n$: number of haplotype samples (indexed $0, \ldots, n{-}1$)
\item $m$: number of variant sites
\item $n_i$: number of global intervals (unique segment boundaries)
\item $S_i$: number of segments for sample $i$; $S = \sum_i S_i$
\item $M_i$: number of mismatches for sample $i$; $M = \sum_i M_i$
\item $d$: depth of the threading dependency tree ($\le \lceil\log_2 n\rceil$ typically)
\item $\bar{S} = S/n$: average segments per sample ($\approx 20$ for 1\,Mb regions)
\item $T$: number of OpenMP threads
\end{itemize}

%% ========================================================================
\subsection{Preparation}

Preparation is a one-time $O(nS\log n_i + Md + nSd)$ computation that
builds all data structures needed for subsequent multiply calls.  It
proceeds in four phases.

\subsubsection{Phase 1: Global Interval Decomposition}

Collect all segment start positions across all $n$ samples into a sorted,
deduplicated set of boundary positions $\{b_0, b_1, \ldots, b_{n_i}\}$.
The $j$-th interval spans sites $[\texttt{ivl\_start}[j],
\texttt{ivl\_end}[j])$, where $\texttt{ivl\_start}[j]$ and
$\texttt{ivl\_end}[j]$ are the first and one-past-last \emph{site indices}
within the interval.  Within each interval, every sample's target is
constant.

\textbf{Auxiliary:} Build the reference genome
$\texttt{ref}[s] \in \{0,1\}$ for sample~0 by marking each mismatch of
sample~0 from the all-zero baseline.  Build $\texttt{site\_to\_ivl}[s]$
mapping each site to its containing interval ($O(m)$ memory, one scan).

\textbf{Cost:} $O(S \log S + m)$.

\subsubsection{Phase 2: Per-Sample Segment-to-Interval Mapping}

For each sample $i$, map its $S_i$ segments to global intervals.
$\texttt{tree\_seg\_first\_ivl}[\texttt{offset}[i] + k]$ stores the first
global interval of sample $i$'s $k$-th segment, computed via binary search
on $\texttt{ivl\_start}$.

This enables $O(\log S_i)$ lookup of ``which target does sample $i$ have
at interval $j$?''\ via the helper function
$\texttt{target\_at\_interval}(i, j)$, which binary-searches
$\texttt{tree\_seg\_first\_ivl}$ for the largest entry $\le j$.

\textbf{Cost:} $O(S \log n_i)$.  \textbf{Memory:} $O(S)$.

\subsubsection{Phase 3: Mismatch Signs via Lazy Chain Tracing}

For each non-self mismatch of each sample, compute the correction sign
$\texttt{mm\_sign}[k] \in \{-1, +1\}$ by querying the target's genotype at
that site.  The genotype is determined by tracing the threading chain from
the target back to sample~0, accumulating XOR parity of mismatches along
the path.  Each query costs $O(d \cdot \log \bar{S})$ (binary search at
each of $d$ chain levels to find the active segment).

Self-referencing mismatches ($\text{target} = i$) are marked with
$\texttt{mm\_sign} = 0$ and handled specially during multiply.

\textbf{Key property:} No genotype buffer is materialized.  The entire
phase uses $O(1)$ working memory beyond the output arrays.

\textbf{Cost:} $O(Md \log \bar{S})$.  \textbf{Memory:} $O(M)$.

\subsubsection{Phase 4: Multiply Support Structures}

This phase builds four data structures used during the interval sweep of
both left and right multiply:

\paragraph{A. Inverted mismatch index.}
Group non-self-ref mismatches by interval for $O(1)$ iteration during the
sweep.  For each interval $j$, the range
$[\texttt{inv\_mm\_offset}[j], \texttt{inv\_mm\_offset}[j{+}1])$ indexes
into flat arrays $\texttt{inv\_mm\_sample}$, $\texttt{inv\_mm\_site}$,
$\texttt{inv\_mm\_sign}$, giving all (sample, site, sign) triples whose
mismatch falls in interval $j$.

Built by scanning $\texttt{site\_to\_ivl}$ for each mismatch, filtering
out self-ref ($\texttt{mm\_sign} = 0$), and counting-sort by interval.
Memory: $O(M)$.

\paragraph{B. Interval change list.}
Record which samples change target at each interval boundary $j \to j{+}1$.
For each change at boundary $j$, store
$(\texttt{sample}, \texttt{old\_tgt}, \texttt{new\_tgt})$ in flat arrays
indexed by $[\texttt{ivl\_change\_offset}[j],
\texttt{ivl\_change\_offset}[j{+}1])$.

Within each group, entries are sorted by \emph{decreasing sample index}.
This ordering is critical for correctness of left multiply boundary updates
(Section~\ref{sec:left-boundary}).

Built by scanning $\texttt{tree\_seg\_first\_ivl}$: at each segment
boundary for sample $i$ at interval $j$, record the old and new targets.
Memory: $O(S)$.

\paragraph{C. Self-ref sample list.}
For each interval $j$, store which samples are self-referencing
($\text{target} = i$) in flat arrays
$\texttt{self\_ref\_sample}[\texttt{self\_ref\_offset}[j] \ldots]$.
Memory: $O(S_\text{self})$ (typically small).

\paragraph{D. Precomputed ancestor chains.}
For each entry in the interval change list, precompute the full ancestor
chain that must be updated in the left multiply's $W$ vector.  Two sets of
chains are stored:

\emph{Detach chains:} For change $c$ at boundary $j$, the detach chain is
the sequence of ancestors of $\texttt{old\_tgt}$ in the tree at interval
$j$: starting from $\texttt{old\_tgt}$, repeatedly look up the target at
interval $j$ until reaching sample~0.  Stored in
$\texttt{detach\_chain\_data}[\texttt{detach\_chain\_offset}[c] \ldots]$.

\emph{Attach chains:} Same, but for $\texttt{new\_tgt}$ in the tree at
interval $j{+}1$.  Stored in $\texttt{attach\_chain\_data}$.

Precomputing these chains during preparation (total cost $O(nSd)$, memory
$O(nSd)$) eliminates the $O(\log \bar{S})$ binary search per ancestor step
that would otherwise be needed during each left multiply call.

\textbf{Total Phase 4 cost:} $O(M + S + nSd)$.

%% ========================================================================
\subsection{Right Multiply: $\mathbf{y} = \mathbf{G}^\top \mathbf{x}$}

Given $\mathbf{x} \in \mathbb{R}^m$, compute $y_i = \sum_s G_{si} \,
x_s$ for each sample $i$.

The algorithm sweeps intervals $j = 0, \ldots, n_i{-}1$, maintaining a
\emph{propagation vector} $\delta[i]$ that accumulates each sample's
dot-product contribution from the current interval.  At the end of each
interval, weights propagate up the threading tree: for $i = n{-}1$ down to
$1$, $\delta[\text{target}_i] \mathrel{+}= \delta[i]$, then
$\delta[i] = 0$.

\begin{algorithm}[h]
\caption{Tree-shuttle right multiply (single-threaded core)}
\begin{algorithmic}[1]
\State Allocate $\delta[0 \ldots n{-}1] \gets 0$,
       $\texttt{mm\_corr}[0 \ldots n{-}1] \gets 0$,
       $y[0 \ldots n{-}1] \gets 0$
\State $\texttt{total\_ref} \gets 0$
\For{$j = 0$ to $n_i - 1$}
    \If{$\texttt{ivl\_start}[j] \ge \texttt{ivl\_end}[j]$}
        \State \textbf{continue} \Comment{skip empty intervals}
    \EndIf
    \State \textbf{// A. Reference contribution}
    \State $\texttt{ref\_sum} \gets \sum_{s=\texttt{ivl\_start}[j]}^{\texttt{ivl\_end}[j]-1} x[s] \cdot \texttt{ref}[s]$
    \State $\texttt{total\_ref} \mathrel{+}= \texttt{ref\_sum}$
    \State $\delta[0] \mathrel{+}= \texttt{ref\_sum}$
    \State \textbf{// B. Mismatch corrections}
    \For{each $(i, s, \sigma)$ in $\texttt{inv\_mm}$ for interval $j$}
        \State $\texttt{mm\_corr}[i] \mathrel{+}= \sigma \cdot x[s]$
    \EndFor
    \State \textbf{// C. Propagation (bottom-up)}
    \For{$i = n - 1$ down to $1$}
        \State $t \gets \texttt{cur\_target}[i]$
        \If{$t = i$} \Comment{self-ref: use carry}
            \State handle via carry-run decomposition
        \Else
            \State $y[i] \mathrel{+}= \texttt{mm\_corr}[i]$;
                   $\texttt{mm\_corr}[i] \gets 0$
            \State $\delta[t] \mathrel{+}= \delta[i]$;
                   $\delta[i] \gets 0$
        \EndIf
    \EndFor
    \State $y[0] \mathrel{+}= \texttt{mm\_corr}[0]$;
           $\texttt{mm\_corr}[0] \gets 0$
    \State \textbf{// D. Update targets at boundary}
    \For{each change $(c, t_\text{old}, t_\text{new})$ at boundary $j \to j{+}1$}
        \State $\texttt{cur\_target}[c] \gets t_\text{new}$
    \EndFor
\EndFor
\State $\forall i$: $y[i] \mathrel{+}= \texttt{total\_ref}$
\end{algorithmic}
\end{algorithm}

\textbf{Key insight (deferred reference sum):}  Because every sample
inherits the reference genome's 1-sites (modified only at mismatches), the
reference contribution $\sum_s x[s] \cdot \texttt{ref}[s]$ is common to
all samples.  Rather than propagating it through $\delta$ at every
interval, we accumulate $\texttt{total\_ref}$ and add it once at the end
($O(n)$ instead of $O(n \cdot n_i)$).

\textbf{Complexity:} $O(n \cdot n_i^\text{nonempty} + M)$ per call, where
$n_i^\text{nonempty} \le n_i$ is the number of intervals containing at
least one site.  The propagation loop (step~C) dominates at
$O(n)$ per non-empty interval.  Memory: $O(n)$.

%% ========================================================================
\subsection{Left Multiply: $\mathbf{y} = \mathbf{G} \mathbf{x}$}
\label{sec:left-multiply}

Given $\mathbf{x} \in \mathbb{R}^n$, compute $y_s = \sum_i G_{si} \,
x_i$ for each site $s$.  The key innovation over the previous $O(n \cdot
n_i)$ weight-matrix approach is maintaining an $O(n)$ \emph{subtree weight
vector} $W[i]$ incrementally across intervals.

\textbf{Invariant:} At interval $j$, $W[i]$ equals the sum of $x_k$ over
all samples $k$ in the subtree rooted at $i$ in the threading tree at
interval $j$.  In particular, $W[0] = \sum_i x_i$ (total weight).

\begin{algorithm}[h]
\caption{Tree-shuttle left multiply (single-threaded core)}
\begin{algorithmic}[1]
\State \textbf{Initialize $W$} for interval 0:
\State $W[i] \gets x[i]$ for all $i$
\For{$i = n - 1$ down to $1$}
    \State $t \gets \texttt{target\_at\_interval}(i, 0)$
    \If{$t \ne i$} $W[t] \mathrel{+}= W[i]$ \EndIf
\EndFor
\State \textbf{Sweep intervals:}
\For{$j = 0$ to $n_i - 1$}
    \State \textbf{// A. Reference emission}
    \For{each site $s$ in interval $j$ where $\texttt{ref}[s] = 1$}
        \State $y[s] \mathrel{+}= W[0]$
    \EndFor
    \State \textbf{// B. Mismatch emission}
    \For{each $(i, s, \sigma)$ in $\texttt{inv\_mm}$ for interval $j$}
        \State $y[s] \mathrel{+}= \sigma \cdot W[i]$
    \EndFor
    \State \textbf{// C. Self-ref carry emission}
    \For{each self-ref sample $i$ in interval $j$}
        \State emit carry-run contribution with weight $W[i]$
    \EndFor
    \State \textbf{// D. Boundary update (Section~\ref{sec:left-boundary})}
    \If{$j + 1 < n_i$}
        \State update $W$ via detach/attach passes
    \EndIf
\EndFor
\end{algorithmic}
\end{algorithm}

\subsubsection{Boundary Update: Detach/Attach}
\label{sec:left-boundary}

When transitioning from interval $j$ to $j{+}1$, some samples change
target.  The $W$ vector must be updated to reflect the new tree topology.
This is done in two passes using the precomputed ancestor chains:

\textbf{Pass 1 --- Detach (decreasing sample order):}
For each changing sample $c$ (processed in decreasing index order), subtract
$W[c]$ from every ancestor along the chain from $\texttt{old\_tgt}(c)$ up
to sample~0 in the interval-$j$ tree:
\begin{algorithmic}[1]
\For{each change $c$ in decreasing sample order}
    \State $\delta_c \gets W[c]$
    \For{each ancestor $a$ in $\texttt{detach\_chain}[c]$}
        \State $W[a] \mathrel{-}= \delta_c$
    \EndFor
\EndFor
\end{algorithmic}

\textbf{Pass 2 --- Attach (arbitrary order):}
For each changing sample $c$, add the \emph{saved} $\delta_c$ to every
ancestor along the chain from $\texttt{new\_tgt}(c)$ up to sample~0 in the
interval-$(j{+}1)$ tree:
\begin{algorithmic}[1]
\For{each change $c$}
    \For{each ancestor $a$ in $\texttt{attach\_chain}[c]$}
        \State $W[a] \mathrel{+}= \delta_c$
    \EndFor
\EndFor
\end{algorithmic}

\textbf{Correctness of decreasing order:}  Processing detachments in
decreasing sample order ensures $W[c]$ is correct before detaching, because
targets always point to lower-indexed samples.  A higher-indexed descendant
of $c$ that also changes at this boundary is detached first, so its weight
is already removed from $W[c]$.

\textbf{Complexity per boundary:} $O(|\text{changes}| \cdot d)$ using
precomputed chains.  Total across all intervals:
$O(\text{total\_changes} \cdot d) = O(S \cdot d)$.

\textbf{Per-call complexity:} $O(m + M + Sd)$.  Memory: $O(n + m)$.

%% ========================================================================
\subsection{Region-Chunked Multithreading}
\label{sec:multithreading}

Both left and right multiply parallelize via \emph{region chunking}: the
$n_i$ intervals are divided into $T$ contiguous chunks, one per thread.
Each thread independently processes its chunk.

\subsubsection{Thread Initialization}

Each thread must establish its starting state at interval $j_\text{start}$:

\textbf{Right multiply:} Initialize $\texttt{cur\_target}[i]$ for all $n$
samples by calling $\texttt{target\_at\_interval}(i, j_\text{start})$,
which performs $O(\log \bar{S})$ binary search per sample.  Total: $O(n
\log \bar{S})$ per thread.  All per-sample accumulators ($\delta$,
$\texttt{mm\_corr}$) start at zero.

\textbf{Left multiply:} Build the full $W$ vector for the starting
interval.  Initialize $W[i] = x[i]$, then propagate bottom-up: for $i =
n{-}1$ down to $1$, look up the target at $j_\text{start}$ and accumulate
$W[\text{target}] \mathrel{+}= W[i]$.  Total: $O(n \log \bar{S})$ per
thread.

\subsubsection{Right Multiply Parallelism}

Each thread produces a local $\texttt{out}[n]$ vector from its chunk.
After the sweep, local vectors are combined via a critical-section
reduction ($O(n)$ per thread, negligible).

\begin{algorithmic}[1]
\State \textbf{parallel} $T$ threads:
\State \quad chunk $[j_\text{start}, j_\text{end}) \gets$ thread's share of $[0, n_i)$
\State \quad initialize $\texttt{cur\_target}$ at $j_\text{start}$
\State \quad $\texttt{local\_out}[n] \gets 0$
\State \quad sweep intervals $j_\text{start}$ to $j_\text{end} - 1$ (Algorithm 5 core)
\State \quad \textbf{critical:} $y[i] \mathrel{+}= \texttt{local\_out}[i]$ for all $i$
\end{algorithmic}

\subsubsection{Left Multiply Parallelism}

Because intervals partition the sites, each thread writes to a disjoint
range of the output array $y[s]$.  \emph{No reduction is needed.}

\begin{algorithmic}[1]
\State \textbf{parallel} $T$ threads:
\State \quad chunk $[j_\text{start}, j_\text{end}) \gets$ thread's share of $[0, n_i)$
\State \quad build $W$ for $j_\text{start}$
\State \quad sweep intervals $j_\text{start}$ to $j_\text{end} - 1$ (Algorithm 6 core)
\State \quad \Comment{writes directly to shared $y[\cdot]$, non-overlapping}
\end{algorithmic}

\subsubsection{Thread Scaling}

Let $C_\text{init}$ be the per-thread initialization cost ($O(n \log
\bar{S})$, fixed) and $C_\text{work}$ be the total sweep work (scales with
region length).  The parallel time is approximately:
\[
t_T \approx C_\text{init} + \frac{C_\text{work}}{T}.
\]
For short regions (5\,Mb), $C_\text{init}$ is significant, limiting
speedup to $\sim$2--3$\times$ at $T = 8$.  For chromosome-length regions
(50+\,Mb), work dominates and near-linear speedup ($\sim$$T \times$) is
achieved.

%% ========================================================================
\subsection{Batch Multiply}

Processing $k$ vectors simultaneously avoids $k$ separate tree traversals.
Both batch functions use the same region-chunked OpenMP strategy as their
single-vector counterparts: each thread handles a contiguous chunk of
intervals, initializes its own state at the chunk boundary, and sweeps
independently.

\textbf{Right multiply batch:}
Thread-local $\delta$, \texttt{mm\_corr}, and \texttt{local\_out} arrays
are widened to $n \times k$.  Inner loops over $k$ elements auto-vectorize
via SIMD.  Each thread accumulates into its own \texttt{local\_out} and
a critical-section reduction combines results.

\textbf{Left multiply batch:}
Each thread builds a $k$-wide $W$ vector ($n \times k$) for its starting
interval and sweeps its chunk.  Precomputed ancestor chains are used for
boundary updates: for each detach/attach step, a length-$k$ vector
subtraction or addition is performed.  Since intervals partition the sites,
threads write to non-overlapping ranges of the output---no reduction needed.

%% ========================================================================
\subsection{Complexity Summary}

\begin{table}[h]
\centering
\begin{tabular}{l c c c}
\toprule
& Prepare & Right multiply & Left multiply \\
\midrule
Time & $O(S\!\log n_i + Md\!\log\bar{S} + nSd)$
     & $O\!\left(\frac{n \cdot n_i + M}{T}\! + nT\!\log\bar{S}\right)$
     & $O\!\left(\frac{m + M + Sd}{T}\! + nT\!\log\bar{S}\right)$ \\[4pt]
Memory & $O(S + M + nSd)$ & $O(nT)$ & $O(nT + m)$ \\
\bottomrule
\end{tabular}
\caption{Complexity of tree-shuttle preparation and per-call multiply.
$T$: number of threads.  The $nT\!\log\bar{S}$ terms are per-thread
initialization costs (target lookup via binary search).  Left multiply
memory is $O(n)$ per thread (the $W$ vector) plus $O(m)$ for the shared
output.}
\label{tab:complexity}
\end{table}

\begin{table}[h]
\centering
\begin{tabular}{l c c c}
\toprule
Method & Prepare & Right per-call & Left per-call \\
\midrule
Dense materialization  & $O(nm)$  & $O(nm)$ & $O(nm)$ \\
Tree shuttle (old)     & $O(Md)$  & $O(n \cdot n_i + M)$ & $O(n \cdot n_i + M)$ \\
Tree shuttle (current) & $O(nSd)$ & $O(n \cdot n_i / T + M)$ & $O((m + M + Sd)/T)$ \\
RLE                    & $O(nm)$  & $O(m + R)$ & $O(m + R)$ \\
\bottomrule
\end{tabular}
\caption{Comparison of multiply methods.  The old tree shuttle used an
$O(n \cdot n_i)$ weight matrix for left multiply, which is infeasible at
biobank scale.  The current version replaces this with an incremental
$O(n)$ weight vector, making left multiply sublinear in $n$.  RLE
preparation requires $O(nm)$ full genotype reconstruction.}
\label{tab:comparison}
\end{table}

\subsection{Memory at Biobank Scale}

\begin{table}[h]
\centering
\begin{tabular}{l r r}
\toprule
Component & 500K array (5\,Mb) & 500K WGS (50\,Mb) \\
\midrule
\multicolumn{3}{c}{\emph{Preparation (one-time)}} \\
Segment mapping ($3S$ ints)  & 240\,MB & 24\,GB \\
Mismatch signs ($M$ bytes)   & 1\,GB   & 10\,GB \\
Inverted mismatch index      & 5\,GB   & 50\,GB \\
Interval change list         & 240\,MB & 24\,GB \\
Ancestor chains              & 1.4\,GB & 144\,GB \\
\midrule
\textbf{Prepare total}       & $\sim$8\,GB & $\sim$250\,GB \\
\midrule
\multicolumn{3}{c}{\emph{Per-call}} \\
Right multiply               & 32\,MB  & 32\,MB \\
Left multiply                & 36\,MB  & 36\,MB \\
Old left multiply ($W$ matrix) & \textbf{200\,GB} & \textbf{2\,TB} \\
\bottomrule
\end{tabular}
\caption{Projected memory for 500K haplotypes.  Assumes $\bar{S} = 100$
segments/sample, $\bar{M} = 2000$ mismatches/sample (array, 5\,Mb) or
$\bar{S} = 10000$, $\bar{M} = 20000$ (WGS, 50\,Mb).  Tree depth $d
\approx 15$.  Per-call memory is $O(nT)$: one $W$ or $\delta$ vector per
thread.  The old left multiply's $O(n \cdot n_i)$ weight matrix is shown
for comparison---it is infeasible at this scale.}
\label{tab:memory}
\end{table}

The critical improvement is that per-call left multiply memory dropped from
$O(n \cdot n_i) \approx 200$\,GB to $O(n) \approx 4$\,MB (at $n =
50000$), enabling biobank-scale left multiply that was previously
impossible with the tree-shuttle method.

%% ========================================================================
\subsection{Empirical Scaling}

Table~\ref{tab:bench} reports benchmarks on \textsc{Threads}-inferred ARGs
from msprime simulations (1\,Mb region, varying $n$, 8 OpenMP threads).

\begin{table}[h]
\centering
\begin{tabular}{r r r r r r}
\toprule
$n$ & $m$ & Prepare (ms) & Left (ms) & Right (ms) & Left mem \\
\midrule
100 & 488 & 0.1 & 0.01 & 0.01 & 0\,MB \\
500 & 1546 & 0.5 & 0.04 & 0.18 & 0\,MB \\
1000 & 2577 & 2 & 0.07 & 0.47 & 0\,MB \\
5000 & 8454 & 8 & 0.22 & 3.8 & 0\,MB \\
10000 & 14310 & 19 & 0.35 & 8.5 & 0\,MB \\
50000 & 50715 & 79 & 1.0 & 32 & 4.8\,MB \\
\bottomrule
\end{tabular}
\caption{Tree-shuttle multiply benchmarks (8 threads, 5\,Mb region).
Left multiply memory is $O(n)$; the 4.8\,MB at $n{=}50000$ is
$8 \times 50000 \times 8 = 3.2$\,MB for 8 thread-local $W$ vectors
plus output buffer.}
\label{tab:bench}
\end{table}

\textbf{Scaling exponents} (fitted from $n = 10000$ to $50000$):
\begin{itemize}
\item Prepare: $\sim n^{1.05}$ (near-linear; dominated by chain precomputation)
\item Left multiply: $\sim n^{0.53}$ (sublinear; mismatches grow slower than $n$)
\item Right multiply: $\sim n^{1.03}$ (linear; $O(n)$ propagation per interval)
\end{itemize}

\textbf{Extrapolation to 500K:}  Left $\approx 7.6$\,ms, right $\approx
327$\,ms (8 threads).  Compared to the old tree-shuttle left multiply at
$n = 50000$ (392\,ms, 1.7\,GB memory), the new version is 381$\times$
faster and uses 362$\times$ less memory.

\textbf{Thread scaling at $n = 50000$:}

\begin{table}[h]
\centering
\begin{tabular}{r r r r r}
\toprule
Threads & Left (ms) & Right (ms) & Left speedup & Right speedup \\
\midrule
1 & 2.31 & 197 & 1.0$\times$ & 1.0$\times$ \\
2 & 1.58 & 102 & 1.5$\times$ & 1.9$\times$ \\
4 & 1.10 & 52 & 2.1$\times$ & 3.8$\times$ \\
8 & 0.92 & 31 & 2.5$\times$ & 6.4$\times$ \\
\bottomrule
\end{tabular}
\caption{Thread scaling at $n = 50000$, 5\,Mb region.  Right multiply
achieves near-linear scaling because the $O(n \cdot n_i)$ sweep work
dominates the $O(n \log \bar{S})$ initialization.  Left multiply is
limited by initialization overhead at this region length; for 50\,Mb
regions, left multiply speedup would approach $\sim$6$\times$ at 8
threads.}
\label{tab:threads}
\end{table}

%% ========================================================================
\section{Implementation Notes}

\subsection{Lazy Chain Tracing}

Mismatch correction signs (Phase~3) and target lookups are computed without
materializing any genotype buffer.  For each query
$\texttt{genotype\_at\_site}(i, s)$, the threading chain is traced from
sample $i$ back to sample~0: at each step, binary-search the current
sample's segments to find the active target, check whether $s$ is a
mismatch for this sample, XOR the parity, then recurse on the target.  The
chain terminates at sample~0, whose genotype is determined by its mismatch
list alone.

Each query costs $O(d \cdot \log \bar{S})$.  The total Phase~3 cost is
$O(Md \log \bar{S})$, which is sublinear in $nm$ for typical parameters.

\subsection{Pending Clade/Branch Deferral}

In visit algorithms, clades and branches are not immediately emitted.
Instead, they are stored in a pending list and emitted when the tree
topology changes.  At emission time, the genomic span is set to
$[\text{start\_bp}, b_\text{next})$, spanning all consecutive intervals
with identical topology.  This reduces the number of callbacks and allows
downstream consumers to process contiguous genomic ranges.

\subsection{Interval Skipping}

Empty intervals (containing zero sites) are skipped during both left and
right multiply sweeps.  In right multiply, the propagation loop ($O(n)$
per interval) is the dominant cost; skipping empty intervals can reduce
the effective interval count by 30--50\% in typical threading instructions.

\subsection{Self-Referencing Segments}

A sample targeting itself ($\text{target}_i = i$) has genotype determined
entirely by its mismatches from a carry state.  In right multiply, the
contribution is computed via prefix-sum lookups on mismatch runs.  In left
multiply, the contribution uses the subtree weight $W[i]$ (which equals
$x[i]$ since a self-ref sample has no descendants in that interval)
multiplied by the carry-run structure.

\begin{thebibliography}{9}
\bibitem{gunnarsson2025}
B.~Gunnarsson, M.~Palacios, P.~P.~Palacios, and P.~F.~Palamara.
\newblock Inferring genome-wide genealogies for large samples.
\newblock \emph{In preparation}, 2025.
\end{thebibliography}

\end{document}
